\documentclass[12pt]{article}

\usepackage{amsmath,amssymb}
\usepackage{tikz}
\usetikzlibrary{positioning,arrows.meta,calc}
\usepackage{geometry}
\usepackage{microtype}
\usepackage{booktabs}

\geometry{margin=1in}

% Help avoid overfull lines in text-heavy notes.
\setlength{\emergencystretch}{3em}

\title{Sparse PCA: Derivation, Intuition, and Visualization}
\author{}
\date{}

\begin{document}

\maketitle

\section{Sparse PCA Objective}

Classical PCA solves
\[
\min_{V} \; \|X - ZV^\top\|_F^2, \qquad \text{s.t. } V^\top V = I.
\]
Here
\begin{itemize}
    \item $X \in \mathbb{R}^{n \times p}$ is the data matrix,
    \item $V \in \mathbb{R}^{p \times k}$ are (dense) loadings,
    \item $Z \in \mathbb{R}^{n \times k}$ are principal component scores.
\end{itemize}

Sparse PCA seeks \emph{sparse} loadings $\beta_i$ that approximate the dense principal components $Z_i$:
\[
\min_{\beta_i} \; \|Z_i - X\beta_i\|_2^2 + \lambda \|\beta_i\|_2^2 + \lambda_1 \|\beta_i\|_1,
\]
where
\begin{itemize}
    \item $Z_i = Xv_i$ is the $i$th dense score vector,
    \item $\beta_i \in \mathbb{R}^p$ is the sparse loading vector,
    \item $\lambda$ and $\lambda_1$ control $\ell_2$ and $\ell_1$ regularization.
\end{itemize}

\section{Matrix Form Derivation (One Component)}

Start with the one-component reconstruction objective:
\[
\min_{v_i} \; \|X - z_i v_i^\top\|_F^2.
\]
Using $\|A\|_F^2 = \mathrm{tr}(A^\top A)$,
\[
\|X - z_i v_i^\top\|_F^2 = \mathrm{tr}\big((X - z_i v_i^\top)^\top (X - z_i v_i^\top)\big).
\]
Expanding gives
\[
\mathrm{tr}(X^\top X)
- 2\,\mathrm{tr}(v_i^\top X^\top z_i)
+ \mathrm{tr}\big(v_i^\top (z_i^\top z_i) v_i\big).
\]
Define $\beta_i := v_i$ and add $\ell_1$/$\ell_2$ penalties:
\[
\min_{\beta_i}
\; \mathrm{tr}\big(\beta_i^\top (z_i^\top z_i) \beta_i\big)
- 2\,\mathrm{tr}(\beta_i^\top X^\top z_i)
+ \lambda\,\mathrm{tr}(\beta_i^\top \beta_i)
+ \lambda_1\|\beta_i\|_1.
\]
Since $z_i^\top z_i$ is a scalar, this is equivalent (up to an additive constant) to
\[
\min_{\beta_i} \; \|Z_i - X\beta_i\|_2^2 + \lambda \|\beta_i\|_2^2 + \lambda_1 \|\beta_i\|_1,
\]
which is a regression problem with sparsity.

\section{Intuition}

\begin{itemize}
    \item Dense loadings $v_i$ capture maximum variance in the score $Z_i$.
    \item Sparse loadings $\beta_i$ approximate $Z_i$ using only a subset of features.
    \item The objective trades off fit (small $\|Z_i - X\beta_i\|_2^2$) and sparsity (small $\|\beta_i\|_1$).
\end{itemize}

\noindent
\textbf{Regression perspective:} Sparse PCA can be seen as regressing the dense score $Z_i$ on the original features $X$ to find sparse weights.

\section{Summary Notes}

\begin{itemize}
    \item Sparse PCA approximates dense PCA components using sparse linear combinations of features.
    \item The regression form connects Sparse PCA to LASSO / elastic net solvers.
    \item The $\ell_1$ penalty induces sparsity; the $\ell_2$ penalty stabilizes the solution.
    \item Each component $\beta_i$ can be solved via component-wise regression.
\end{itemize}

\section{Regression-type PCA}

We start with a connection between PCA and regression. Suppose $X \in \mathbb{R}^{n \times p}$ is the data matrix, and we want to derive the first principal component using a \emph{regression-type formulation}.

\subsection{Theorem 2: Leading Principal Component via Regression}
\textbf{Theorem 2.} For any $\lambda > 0$, define
\[
(\hat{\alpha}, \hat{\beta}) = \arg \min_{\alpha, \beta} \sum_{i=1}^n \| x_i - \alpha \beta^\top x_i \|_2^2 + \lambda \|\beta\|_2^2
\]
subject to 
\[
\|\alpha\|_2 = 1.
\]

\textbf{Result:} $\hat{\beta} \propto V_1$, where $V_1$ is the first principal component of $X$.  

\textbf{Interpretation:}
\begin{itemize}
    \item $\beta$ is the loading vector (shared across all rows), corresponding to the principal direction in feature space.
    \item $\alpha$ is the scaling vector used to reconstruct $x_i$ from its projection $\beta^\top x_i$.
    \item Each row of the reconstruction: $\hat{x}_i = \alpha (\beta^\top x_i)$, a rank-1 approximation of $x_i$.
\end{itemize}

\subsection{Theorem 3: Multiple Principal Components}

Let us now consider the first $k$ principal components. Define
\[
A = [\alpha_1, \dots, \alpha_k] \in \mathbb{R}^{p \times k}, \quad
B = [\beta_1, \dots, \beta_k] \in \mathbb{R}^{p \times k}.
\]

\textbf{Theorem 3.} For any $\lambda > 0$, define
\[
(A, B) = \arg \min_{A, B} \sum_{i=1}^n \| x_i - A B^\top x_i \|_2^2 + \lambda \sum_{j=1}^k \|\beta_j\|_2^2
\]
subject to
\[
A^\top A = I_{k \times k}.
\]

\textbf{Result:} $\hat{\beta}_j \propto V_j$ for $j = 1, \dots, k$, i.e., solving this regression-type criterion recovers all $k$ principal components simultaneously.

\subsection{Matrix Form of the Objective}

The row-wise sum of squared errors can be written compactly as a Frobenius norm:
\[
\sum_{i=1}^n \| x_i - A B^\top x_i \|_2^2 = \| X - X B A^\top \|_F^2
\]

\subsection{Decomposition using Orthogonal Complement}

Let $A^\perp \in \mathbb{R}^{p \times (p-k)}$ be any orthonormal complement of $A$, so that $[A \quad A^\perp]$ is a $p \times p$ orthonormal matrix. Then we can decompose the error:
\[
\| X - X B A^\top \|_F^2 = \| X A^\perp \|_F^2 + \| X A - X B \|_F^2
= \| X A^\perp \|_F^2 + \sum_{j=1}^k \| X \alpha_j - X \beta_j \|_2^2
\]

\textbf{Key points:}
\begin{itemize}
    \item $\|X A^\perp\|_F^2$ is the part of $X$ outside the column space of $A$ and cannot be reduced by $B$.
    \item Once $A$ is fixed, minimizing over $B$ reduces to $k$ independent ridge regressions:
    \[
    \hat{\beta}_j = \arg \min_{\beta_j} \| X \alpha_j - X \beta_j \|_2^2 + \lambda \|\beta_j\|_2^2, \quad j=1,\dots,k.
    \]
\end{itemize}

\subsection{SPCA: Introducing Sparsity via Lasso}

To obtain sparse loadings, we add a Lasso penalty to the criterion:
\[
(A, B) = \arg \min_{A,B} \sum_{i=1}^n \| x_i - A B^\top x_i \|_2^2 + \lambda \sum_{j=1}^k \|\beta_j\|_2^2 + \sum_{j=1}^k \lambda_{1,j} \|\beta_j\|_1
\]
subject to $A^\top A = I_{k \times k}$.  

\textbf{Notes:}
\begin{itemize}
    \item $\lambda$ stabilizes the solution and ensures exact PCA when $p>n$ and $\lambda_{1,j}=0$.
    \item $\lambda_{1,j}$ controls sparsity for each principal component individually.
    \item This criterion is called the \textbf{SPCA criterion}.
\end{itemize}

\section{Intuition Summary}

\begin{itemize}
    \item Classical PCA: dense loading vectors $V_j$, perfect reconstruction along principal directions.
    \item Regression-type PCA: recast PCA as a \emph{regression problem}, $\beta_j$ as regression coefficients, $\alpha_j$ as scaling vectors.
    \item Orthogonal decomposition: separates the reconstruction error into the part inside the column space of $A$ (optimizable) and the orthogonal part (irreducible).
    \item SPCA: adding Lasso penalty produces sparse $\beta_j$ while approximating classical PCs.
\end{itemize}

\appendix

\section{Additional Notes: Convex Relaxations and Online Sparse PCA}

\subsection{Background and Motivation}

Sparse Principal Component Analysis (Sparse PCA) seeks a low-dimensional principal
subspace while enforcing sparsity in the loading vectors. Relative to classical
PCA, sparsity improves interpretability and can improve statistical performance in
high-dimensional settings where the leading eigenvectors are believed to be sparse.

The difficulty is computational: the most direct formulations impose $\ell_0$
constraints (or related combinatorial constraints) on the loadings, leading to
nonconvex and generally NP-hard optimization problems. As a result, much of the
literature studies convex relaxations and efficient iterative heuristics that trade
statistical accuracy for computational tractability.

\subsection{Baseline: Fantope Projection and Selection}

One of the most prominent convex relaxations for Sparse PCA is the
\emph{Fantope Projection and Selection} (FPS) method.
Instead of directly estimating sparse eigenvectors, FPS optimizes over a matrix
variable constrained to lie in the Fantope, defined as the convex hull of rank-$d$
projection matrices.

Although FPS enjoys strong statistical guarantees and convexity, it is
computationally expensive. In particular, each iteration typically requires
eigen-decomposition and projection onto the Fantope, which becomes prohibitively
costly in high-dimensional settings. As a result, FPS is often impractical for
large-scale or streaming data scenarios.

\subsection{Limitations of Existing Online Sparse PCA Methods}

While online PCA is well understood, extending Sparse PCA to the online setting
poses significant challenges. Existing online Sparse PCA methods are often based on
heuristics or nonconvex updates, and typically lack either statistical guarantees,
computational efficiency, or both.

In fact, designing an online Sparse PCA algorithm that is simultaneously
\begin{itemize}
    \item statistically consistent,
    \item computationally efficient,
    \item globally convergent,
\end{itemize}
under a general covariance model remains an open problem in the literature.

\subsection{Quick Review: PCA and Sparse PCA}

Given centered data $X \in \mathbb{R}^{n \times p}$, the sample covariance is
\[
\Sigma = \frac{1}{n} X^\top X.
\]

\paragraph{Variance maximization view}
Single-component PCA can be formulated as
\[
\max_{v \in \mathbb{R}^p} \; v^\top \Sigma v \quad \text{s.t. } \|v\|_2 = 1.
\]

\paragraph{Regression view}
PCA can also be written as
\[
\min_{Z,V} \; \|X - ZV^\top\|_F^2 \quad \text{s.t. } V^\top V = I_d.
\]

\subsection{Projection Matrix Formulation and Convex Relaxation}

For a unit vector $v$, define the rank-1 projector $P := vv^\top$. Then
$v^\top \Sigma v = \mathrm{tr}(\Sigma P)$.
Writing PCA in terms of $P$ makes the objective linear, while the constraints
(rank/idempotence) are nonconvex.

Relaxing $P^2=P$ and $\mathrm{rank}(P)=d$ to $0 \preceq P \preceq I$ while keeping
$\mathrm{tr}(P)=d$ yields the Fantope, i.e., the convex hull of rank-$d$ projectors.

\end{document}
